% Part: normal-modal-logic
% Chapter: filtrations
% Section: S5-fmp

\documentclass[../../../include/open-logic-section]{subfiles}

\begin{document}

\olfileid{nml}{fil}{fmp}

\olsection{للمنطقين \Log{K} و\Log{S5} خاصية النموذج المنتهي}

\begin{defn}
  للنظام~$\Sigma$ من المنطق الجهوي \emph{خاصية النموذج المنتهي} إذا
  كان صدق !!a{formula}~$!A$ عند عالم في نموذج لـ~$\Sigma$ يقتضي
  صدق~$!A$ عند عالم في نموذج \emph{منتهٍ} لـ~$\Sigma$.
\end{defn}

\begin{prop}\ollabel{prop:K-fmp}
  يملك~\Log{K} خاصية النموذج المنتهي.
\end{prop}

\begin{proof}
  إن~\Log{K} مجموعة ال!!{formula}s الصالحة، فكل نموذج نموذج
  لـ~\Log{K}. وتفيد~\olref[fil]{thm:filtrations} أنه متى كان
  $\mSat{{\OLId{latin-looped}{M}}}{!A}[{\OLId{latin-ordinary}{w}}]$ كان~$\mSat{{\OLId{latin-looped}{M}}^*}{!A}[{[{\OLId{latin-ordinary}{w}}]}]$ في كل ترشيح
  لـ~$\mModel{M}$ عبر مجموعة~${\OLId{greek-variable}{Gamma}}$ من ال!!{formula}s الجزئية
  لـ~$!A$. ولكل !!{formula} عدد منتهٍ من ال!!{formula}s الجزئية،
  فتتناهى~${\OLId{greek-variable}{Gamma}}$. ومن~\olref[fin]{prop:filt-are-finite} نحصل على
  $\card{{\OLId{latin-looped}{W}}^*} \le {\OLMathNumeral{2}}^{\OLId{latin-ordinary}{n}}$، حيث~${\OLId{latin-ordinary}{n}}$ عدد ال!!{formula}s في~${\OLId{greek-variable}{Gamma}}$.
  ثم إن~\Log{K} لا يشترط في النماذج قيدًا آخر، فيكون~$\mModel{M^*}$
  نموذجًا لـ~\Log{K}، وهو المطلوب.
\end{proof}

إذا أردنا إثبات خاصية النموذج المنتهي لمنطق~\Log{L} بالترشيح، فلا بد
أن يبقى ترشيح نموذج~\Log{L} نموذجًا لـ~\Log{L}. وهذا يحتاج في كثير من
الأحيان إلى برهان غير يسير؛ إذ ليس كل ترشيح نموذجًا لـ~\Log{L}.
أما النماذج الكلية فإن ترشيحاتها تبقى كلية، كما نبين الآن.

\begin{prop}\ollabel{prop:univ-fin}
  لتكن~$\mClass{U}$ فئة النماذج الكلية (انظر
  \olref[frd][es5]{prop:S5=univ})، ولتكن~$\mClass{U}_\mathrm{Fin}$
  فئة جميع النماذج الكلية المنتهية. عندئذ تكون أي !!{formula}~$!A$
  صالحة في~$\mClass{U}$ إذا وفقط إذا كانت صالحة في
  $\mClass{U}_\mathrm{Fin}$.
\end{prop}

\begin{proof}
  أما الاستلزام الأول فواضح، لأن النموذج الكلي المنتهي من جملة النماذج
  الكلية. وأما عكسه فنثبته بعكس النقيض. لنفرض أن~$!A$ تكذب عند عالم
  ${\OLId{latin-ordinary}{w}}$ في نموذج كلي~$\mModel{M}$، ولنأخذ~${\OLId{greek-variable}{Gamma}}$ مؤلفة من~$!A$
  وصيغها الجزئية كلها؛ فهذه~${\OLId{greek-variable}{Gamma}}$ منتهية. نأخذ ترشيحًا
  $\mModel{M^*}$ لـ~$\mModel{M}$ عبرها، فيكون~$\mModel{M^*}$ منتهيًا
  بمقتضى~\olref[fin]{prop:filt-are-finite}. ثم تفيد
  \olref[fil]{thm:filtrations} أن~$!A$ تكذب عند~$[{\OLId{latin-ordinary}{w}}]$ في
  $\mModel{M^*}$. ولا ينقص إلا إثبات أن~$\mModel{M^*}$ كلي.
  فمهما أخذنا~${\OLId{latin-ordinary}{u}}$ و${\OLId{latin-ordinary}{v}}$ كان~$Ruv$ بالفرض، فتوجب
  \olref[fil]{defn:filtration}\olref[fil]{defn:filtration-R}
  العلاقة~${\OLId{latin-looped}{R}}^*[{\OLId{latin-ordinary}{u}}][{\OLId{latin-ordinary}{v}}]$ أيضًا.
\end{proof}

\begin{cor}\ollabel{cor:S5fmp}
  يملك~\Log{S5} خاصية النموذج المنتهي.
\end{cor}

\begin{proof}
  تقضي~\olref[frd][es5]{prop:S5=univ} بأن~$!A$ إذا صدقت عند عالم
  في نموذج انعكاسي إقليدي صدقت عند عالم في نموذج كلي. وبمقتضى
  \olref{prop:univ-fin} يمكن اختيار هذا النموذج الكلي منتهيًا: نرشحه
  عبر مجموعة ال!!{formula}s الجزئية لـ~$!A$. ولما كان كل نموذج كلي
  انعكاسيًا وإقليديًا، حصل صدق~$!A$ عند عالم في نموذج انعكاسي إقليدي
  منتهٍ، وتم البرهان.
\end{proof}

\begin{prob}
  أثبت أن كل ترشيح لنموذج متسلسل أو انعكاسي يكون هو أيضًا متسلسلًا أو
  انعكاسيًا، على الترتيب.
\end{prob}

\begin{prob}
  أوجد ترشيحًا غير متناظر (أو غير متعدٍ، أو غير إقليدي) لنموذج متناظر
  (أو متعدٍ، أو إقليدي).
\end{prob}

\end{document}
